\tcbuselibrary{theorems, skins, breakable}

% ---------------------------------------------------------------------------
% Shared counter (theorems and lemmas share the same counter; definitions,
% examples, notes and remarks each have their own, all subordinate to section)
% ---------------------------------------------------------------------------
% Single unified counter for all numbered environments, reset by section.
% Theorems, lemmas, corollaries, definitions, examples and remarks all
% draw from this one sequence, so numbering is global within each section.
\newcounter{thmcounter}[section]
\renewcommand{\thethmcounter}{\thesection.\arabic{thmcounter}}

% ---------------------------------------------------------------------------
% Internal helper: \envtitlesimple{AccentColor}{Label}{Counter}
%   Typesets the title for note/remark: no Lean name, no horizontal rule.
%   The label is rendered in the accent color for visual identity.
% ---------------------------------------------------------------------------
\newcommand{\envtitlesimple}[3]{%
  % #1 = accent color, #2 = label word, #3 = counter macro
  \noindent{\color{#1}\textbf{#2~#3.}}\quad
}

% ---------------------------------------------------------------------------
% Base tcolorbox style shared by all environments
% ---------------------------------------------------------------------------
\tcbset{
  envbase/.style={
    enhanced,
    breakable,
    rounded corners,
    boxrule=0pt,          % no outer border
    frame hidden,
    colback=white,
    left=6pt, right=6pt, top=2pt, bottom=2pt,
    before upper={\setlength{\parskip}{4pt}},
    fontupper=\normalsize,
  },
  % Extra style for environments with a coloured title banner
  envtitled/.style={
    titlerule=0pt,        % no separator line between title and body
    toptitle=1pt,
    bottomtitle=1pt,
  }
}

% ---------------------------------------------------------------------------
% THEOREM  (accent: blue)
% ---------------------------------------------------------------------------
\NewTColorBox{theorem}{ O{} O{} }{%
  envbase,
  envtitled,
  colback=pink!2,
  colbacktitle=pink!70,
  coltitle=pink!2,
  title={%
    \stepcounter{thmcounter}%
    \textbf{Theorem~\thethmcounter}%
    \if\relax\detokenize{#1}\relax\else~\textup{(#1)}\fi
    \hfill
    \if\relax\detokenize{#2}\relax\else{\small\texttt{#2}}\fi
  },
  borderline west={2pt}{0pt}{pink!70},
}

\NewTColorBox{corollary}{ O{} O{} }{%
  envbase,
  envtitled,
  colback=pink!2,
  colbacktitle=pink!70,
  coltitle=pink!2,
  title={%
    \stepcounter{thmcounter}%
    \textbf{Corollary~\thethmcounter}%
    \if\relax\detokenize{#1}\relax\else~\textup{(#1)}\fi
    \hfill
    \if\relax\detokenize{#2}\relax\else{\small\texttt{#2}}\fi
  },
  borderline west={2pt}{0pt}{pink!70},
}

% ---------------------------------------------------------------------------
% LEMMA  (accent: green)
% ---------------------------------------------------------------------------
\NewTColorBox{lemma}{ O{} O{} }{%
  envbase,
  envtitled,
  colback=pink!2,
  colbacktitle=pink!70,
  coltitle=pink!2,
  title={%
    \stepcounter{thmcounter}%
    \textbf{Lemma~\thethmcounter}%
    \if\relax\detokenize{#1}\relax\else~\textup{(#1)}\fi
    \hfill
    \if\relax\detokenize{#2}\relax\else{\small\texttt{#2}}\fi
  },
  borderline west={2pt}{0pt}{pink!70},
}

% ---------------------------------------------------------------------------
% DEFINITION  (accent: orange)
% ---------------------------------------------------------------------------
\NewTColorBox{definition}{ O{} O{} }{%
  envbase,
  envtitled,
  colback=orange!2,
  colbacktitle=orange!70,
  coltitle=orange!2,
  title={%
    \stepcounter{thmcounter}%
    \textbf{Definition~\thethmcounter}%
    \if\relax\detokenize{#1}\relax\else~\textup{(#1)}\fi
    \hfill
    \if\relax\detokenize{#2}\relax\else{\small\texttt{#2}}\fi
  },
  borderline west={2pt}{0pt}{orange!70},
}

% ---------------------------------------------------------------------------
% EXAMPLE  (accent: pink)
% ---------------------------------------------------------------------------
\NewTColorBox{example}{ O{} O{} }{%
  envbase,
  envtitled,
  colback=green!2,
  colbacktitle=green!60,
  coltitle=green!2,
  title={%
    \stepcounter{thmcounter}%
    \textbf{Example~\thethmcounter}%
    \if\relax\detokenize{#1}\relax\else~\textup{(#1)}\fi
    \hfill
    \if\relax\detokenize{#2}\relax\else{\small\texttt{#2}}\fi
  },
  borderline west={2pt}{0pt}{green!60},
}

% ---------------------------------------------------------------------------
% NOTE  (accent: dark-blue) — no arguments
% ---------------------------------------------------------------------------
\NewTColorBox{note}{ }{%
  envbase,
  colback=beige!20,
  before upper={\setlength{\parskip}{4pt}\textcolor{beige!60!black}{\textbf{Note.}}\enspace},
  borderline west={2pt}{0pt}{beige!98!black},
}

% ---------------------------------------------------------------------------
% REMARK  (accent: dark-blue) — no arguments
% ---------------------------------------------------------------------------
\NewTColorBox{remark}{ }{%
  envbase,
  colback=blue!5,
  before upper={\stepcounter{thmcounter}\setlength{\parskip}{4pt}\textcolor{blue}{\textbf{Remark~\thethmcounter.}}\enspace},
  borderline west={2pt}{0pt}{blue!70},
}

% ---------------------------------------------------------------------------
% RULE  (accent: dark-blue) — no arguments, to title
% ---------------------------------------------------------------------------
\NewTColorBox{grayenv}{ }{%
  envbase,
  envtitled,
  colback=gray!5,
  title={},
}

% ---------------------------------------------------------------------------
% FACT — no title, no accent colour, just a vertical rule on the left
% ---------------------------------------------------------------------------
\NewTColorBox{fact}{ }{%
  envbase,
  notitle,
  colback=blue!7,
  before upper={\setlength{\parskip}{4pt}\textcolor{blue}{\textbf{Fact.}}\enspace},
  borderline west={2pt}{0pt}{blue!70},
}

% ---------------------------------------------------------------------------
% PROOF — "Proof." inline, QED square flush-right on the last line
% ---------------------------------------------------------------------------
\RenewTColorBox{proof}{ }{%
  envbase,
  notitle,
  colback=blue!2,
  before upper={\setlength{\parskip}{4pt}\textit{Proof.}\enspace},
  after upper={\unskip\nobreak\hspace*{1em}\hfill\ensuremath{\square}},
  borderline west={2pt}{0pt}{blue!70},
}

% ---------------------------------------------------------------------------
% SUBPROOF — "Proof." inline, QED square flush-right on the last line.
% ---------------------------------------------------------------------------
\NewTColorBox{subproof}{ }{%
  envbase,
  notitle,
  colback=blue!7,
  before upper={\setlength{\parskip}{4pt}\textcolor{blue}{\textit{Proof.}}\enspace},
  after upper={\unskip\nobreak\hspace*{1em}\hfill\ensuremath{\square}},
  borderline west={2pt}{0pt}{blue!70},
}

% Wrap proof/subproof so that \footnote inside them goes to the page bottom.
% \makesavenoteenv injects \if@endpe bookkeeping that adds spurious vertical
% space, so we save the tcolorbox-defined commands directly and rewrap them
% with plain \renewenvironment — \savenotes intercepts footnotes before the
% box is built; \spewnotes flushes them to \footins after it is placed.
% NB: \makeatletter/\makeatother are required because the internal names
%     contain @, which has catcode 12 outside of package/class code.
\makeatletter
\let\tcb@proof\proof         \let\endtcb@proof\endproof
\let\tcb@subproof\subproof   \let\endtcb@subproof\endsubproof
\renewenvironment{proof}   {\savenotes\tcb@proof}   {\endtcb@proof\spewnotes}
\renewenvironment{subproof}{\savenotes\tcb@subproof}{\endtcb@subproof\spewnotes}
\makeatother